\chapter{Cómo se construyó el proyecto}
\label{cap:construccion}

\section{Por qué aparece Claude en el proyecto}
\textbf{Lo que pasó, tal cual.} El código, las notas y los documentos se construyeron trabajando con \textbf{Claude
Code}, un asistente de programación con inteligencia artificial (de la empresa Anthropic), en sesiones de trabajo con
Brandon. Por eso aparece en varios lugares del repositorio:

\begin{center}
\begin{tabularx}{\linewidth}{lY}
\encabezado \h{Dónde aparece} & \h{Qué es} \\
\texttt{Co-Authored-By: Claude} en cada commit & Git registra que el asistente participó en ese cambio. \\
\codigo{CLAUDE.md} en la raíz & Las \textbf{reglas que Brandon le dio} al asistente: no inventar datos, ir una fase a la vez, que Brandon decide, documentar cada ecuación. \\
La carpeta \codigo{.claude/} & La configuración del asistente (permisos y ganchos). \\
\codigo{graphify-out/} & Un mapa del proyecto (grafo de conocimiento) que el asistente consulta para no perder el contexto. \\
``Claude Design'' en la nota 07 & La herramienta con la que se rediseñó la página. \\
\bottomrule
\end{tabularx}
\end{center}

\textbf{Cómo se trabajó, y cómo se puede demostrar con los archivos.} El protocolo está escrito en \codigo{CLAUDE.md}
(sección 2) y se siguió en todas las fases:
\begin{enumerate}
\item El asistente explicaba qué había que decidir y \textbf{con qué cifra real}.
\item Presentaba \textbf{2 o 3 opciones} con pros y contras.
\item \textbf{Brandon elegía.} Hay más de 40 registros de decisión en \codigo{OBJETIVO.md} (sección A.8) y 25 notas en
  \codigo{docs/decisiones/}, cada una con las palabras textuales de Brandon, las opciones descartadas y la evidencia.
\item El asistente programaba \textbf{una pieza pequeña} (una función o un notebook corto) comentada en español.
\item La pieza se corría, se revisaba con datos reales y quedaba con pruebas automáticas.
\end{enumerate}

\begin{teprofe}
Lo honesto es declararlo antes de que pregunten. Una forma de decirlo: \emph{``Usé un asistente de programación con
inteligencia artificial como herramienta. Las decisiones del proyecto son mías y están documentadas con mis palabras; el
código lo revisé, cada cifra se puede rastrear a su fuente oficial y puedo explicar cada método.''} Conviene revisar
antes la política del profesor y de la UNRC sobre el uso de inteligencia artificial. Lo que hace defendible el proyecto
no es esconder la herramienta, sino poder explicar cada pieza: para eso es esta guía.
\end{teprofe}

\section{Por qué hay tantos archivos de código}
\textbf{La regla:} una pieza = un archivo. Cada archivo hace una sola cosa y empieza con un encabezado obligatorio que
dice qué hace, por qué así (con la alternativa descartada), de qué datos sale y a qué decisión de la campaña alimenta.
Así cada pieza se puede probar sola, explicar sola y reemplazar sin romper las demás. Si todo estuviera en un solo
notebook, no se podría probar por partes ni reutilizar una función en la página.

El paquete \codigo{backend/torre/} tiene \textbf{56 archivos} en 7 carpetas:

\begin{center}
\begin{tabularx}{\linewidth}{lrYl}
\encabezado \h{Carpeta} & \h{Archivos} & \h{Para qué sirve} & \h{Fases} \\
\codigo{base/} & 23 & Descargar las fuentes (Bronze), limpiarlas con Spark (Silver), almacén DuckDB y diccionario & 0, 1, 2 \\
\codigo{radar/} & 7 & Planteamiento, criterios de las regiones, índice de presión, clasificador, Markov, clustering & 3, 4 \\
\codigo{pronostico/} & 7 & Series de tiempo, forma del año, modelos, rango del 90\,\%, escenarios, calendario & 5 \\
\codigo{campana/} & 9 & Presupuesto (IO), minería de texto, marca, lugares (NLP), fotos, idiomas & 6, 8 \\
\codigo{envivo/} & 3 & Señales semanales, motor de reglas, reproducción con Spark Streaming & 7 \\
\codigo{api/} & 2 & Datos de la página y servidor FastAPI & 9 \\
\codigo{documento/} & 5 & Gráficas, PDF, capturas y auditoría de la página, y la carpeta Entregable & 10, 11 \\
\bottomrule
\end{tabularx}
\end{center}

\subsection{El mapa completo, archivo por archivo}
\small
\begin{longtable}{P{0.33\linewidth}P{0.6\linewidth}}
\rowcolor{guinda}\h{Archivo} & \h{Qué hace} \\ \endhead
\multicolumn{2}{l}{\textbf{\color{guinda}base/ — datos crudos y limpios}}\\
\codigo{entorno.py} & Apunta Java 17 y Hadoop solo dentro del proceso y crea la sesión de Spark (\codigo{local[*]}, 4 GB). \\
\codigo{manifiesto.py} & Registra cada archivo crudo con su huella SHA-256, tamaño, URL, fecha y número de registros. \\
\codigo{ingesta_siturq.py} & Descarga los indicadores de SITUR-Q por su API pública. \\
\codigo{ingesta_datatur.py} & Descarga de DataTur la ocupación semanal y mensual, nacionalidades, INAH, AFAC, cruceros y el Compendio. \\
\codigo{ingesta_abiertas.py} & Descarga DENUE, Censo, clima, huracanes, FRED, Rest-Mex, geografía y ENDUTIH. \\
\codigo{ingesta_benchmarks.py} & Abre con un navegador automatizado las tablas de costos de anuncios y guarda captura como evidencia. \\
\codigo{evidencia_sargazo.py} & Guarda como evidencia las publicaciones públicas sobre sargazo en la bahía. \\
\codigo{ingesta_fotos.py} & Fotos con licencia libre de Wikimedia Commons. \\
\codigo{ubicaciones.py} & Comprueba que cada pueblo y zona esté en Quintana Roo y en su municipio. \\
\multicolumn{2}{l}{\textit{Limpieza, un archivo por fuente (12 archivos \codigo{silver_*.py}):}}\\
\codigo{silver_siturq.py} & SITUR-Q en formato largo; Tren Maya sumado por estación; ceros imposibles = hueco. \\
\codigo{silver_datatur_ocupacion.py} & 166 Excel de ocupación en una tabla; gana la versión más reciente de cada semana. \\
\codigo{silver_denue.py} & 6.1 millones de negocios con su categoría turística (PySpark). \\
\codigo{silver_inah.py} & Visitantes por zona y mes; quita un bloque duplicado; marca el papel de cada zona. \\
\codigo{silver_iter.py} & Censo 2020 por localidad; datos reservados por el INEGI quedan vacíos. \\
\codigo{silver_clima.py} & Clima diario 1950–2026 y horario 2019–2026, con hora local. \\
\codigo{silver_huracanes.py} & HURDAT2 con la distancia de cada punto a Chetumal y la regla de tormenta. \\
\codigo{silver_fred.py} & Pesos por dólar diario y mensual sin rellenar feriados; inflación de EE.\,UU. \\
\codigo{silver_nacionalidad.py} & Extranjeros por aeropuerto, país y sexo. \\
\codigo{silver_afac.py} & Pasajeros por aerolínea; resuelve 300 llaves repetidas. \\
\codigo{silver_cruceros.py} & Cruceros por puerto y conciliación con SITUR-Q. \\
\codigo{silver_restmex.py} & 207,873 reseñas; Spark lee el CSV directo. \\
\codigo{almacen.py} & Almacén DuckDB: vistas de Silver y Gold y el modelo estrella. \\
\codigo{diccionario.py} & Genera el diccionario de datos desde el almacén. \\
\multicolumn{2}{l}{\textbf{\color{guinda}radar/ — dónde hay presión}}\\
\codigo{planteamiento.py} & Variables, actores, cuotas y HHI (Fase 3). \\
\codigo{criterios.py} & Tabla de criterios de las regiones. \\
\codigo{panel.py} & Panel lugar $\times$ mes con todas las variables de presión. \\
\codigo{indice.py} & Índice de presión, estados y sensibilidad. \\
\codigo{prediccion.py} & Índice comparable, tres clasificadores con origen móvil, sesgo y predicción. \\
\codigo{markov.py} & Cadena de Markov semanal del norte y su validación. \\
\codigo{clustering.py} & Agrupamiento Ward de 55 centros del país. \\
\multicolumn{2}{l}{\textbf{\color{guinda}pronostico/ — cuándo conviene ir}}\\
\codigo{series.py} & Series a pronosticar y meses que no entrenan. \\
\codigo{forma.py} & Forma del año, fuerza de la temporada y segunda opinión con STL. \\
\codigo{modelos.py} & Cinco modelos y el origen móvil. \\
\codigo{intervalos.py} & Rango del 90\,\% (conformal) y su cobertura real. \\
\codigo{seleccion.py} & Elección del modelo y pronóstico final de 12 meses. \\
\codigo{escenarios.py} & Poisson de tormentas, Monte Carlo y sensibilidad. \\
\codigo{calendario.py} & Temporada alta y recomendaciones del planeador. \\
\multicolumn{2}{l}{\textbf{\color{guinda}campana/ — la campaña}}\\
\codigo{presupuesto.py} & Modelo de optimización de dos etapas (PuLP/CBC). \\
\codigo{texto.py} & Minería de las reseñas: temas, Apriori y palabras. \\
\codigo{marca.py} & Personas, marca, anuncios, medios, calendario e indicadores. \\
\codigo{lugares.py} & NLP por léxico sobre el DENUE: qué hacer, comer y dormir. \\
\codigo{vitrina.py}, \codigo{estrellas.py} & Experiencias, rutas y estrellas oficiales de hoteles. \\
\codigo{fotos_lugares.py}, \codigo{aportes.py}, \codigo{idiomas.py} & Fotos comprobadas, aportes del equipo y los 10 idiomas. \\
\multicolumn{2}{l}{\textbf{\color{guinda}envivo/ — la Torre}}\\
\codigo{senales.py} & Las cuatro señales semanales (con Isolation Forest y ventana deslizante en Spark). \\
\codigo{motor.py} & Reglas de pausa y de ``¿Ibas al norte?''. \\
\codigo{torre.py} & La reproducción con Spark Structured Streaming. \\
\multicolumn{2}{l}{\textbf{\color{guinda}api/ y documento/}}\\
\codigo{datos_pagina.py}, \codigo{servidor.py} & Datos de la página y servidor FastAPI. \\
\codigo{figuras.py}, \codigo{pdf.py}, \codigo{capturas.py}, \codigo{auditoria.py} & Gráficas, PDF, capturas de la página y auditoría de accesibilidad. \\
\codigo{entregable.py} & Arma la carpeta \codigo{Entregable/}: las guías, los cuatro notebooks autosuficientes (con el código de estos módulos dentro), los datos y la página. \\
\end{longtable}
\normalsize

\section{Notebooks, pruebas y reproducibilidad}
\textbf{Notebooks.} Hay seis notebooks narrados (\codigo{01_planteamiento}, \codigo{02_radar}, \codigo{03_pronostico},
\codigo{05_optimizacion}, \codigo{06_torre_en_vivo} y \codigo{07_campana}; el 04 del plan quedó dentro del 03). No
contienen la lógica: \textbf{llaman a las funciones del paquete} y cuentan la historia con texto, tablas y gráficas. Cada
uno se construye y se ejecuta completo con su script \codigo{_construir_*.py}; si una celda falla, el notebook no se
guarda.

\textbf{Los cuatro notebooks de la entrega.} Para los profesores, la carpeta \codigo{Entregable/} junta esos seis en
\textbf{cuatro notebooks autosuficientes}: (1) los datos y la arquitectura, (2) dónde y cuándo (planteamiento, Radar y
Pronóstico), (3) el presupuesto y la Torre en vivo, y (4) la campaña. No dependen de archivos \codigo{.py}: cada celda que
empieza con \codigo{\%\%modulo} trae completo un archivo del paquete, y al correrla queda registrada como ese módulo. Así
el notebook usa exactamente el mismo código y da las mismas cifras (por ejemplo, 956.8 visitantes y 239 lotes de Spark).
Se generan con \codigo{torre.documento.entregable}, que además los ejecuta de principio a fin dentro de la carpeta.

\textbf{Pruebas automáticas (279).} Hay dos tipos:
\begin{itemize}
\item \textbf{De contrato}: revisan la forma de los datos (no hay duplicados, las columnas existen, una tabla tiene las
  semanas que debe).
\item \textbf{De realidad}: revisan una cifra conocida. Por ejemplo, que Tulum haya tenido 1,031,443 visitantes en 2025,
  que la tasa de tormentas de agosto sea 13.9\,\% o que el modelo de presupuesto dé 956.8 visitantes.
\end{itemize}
Las pruebas \codigo{test_documentos.py} y \codigo{test_guias.py} recalculan las cifras que aparecen en los documentos y
en esta guía.

\textbf{Reproducibilidad.} Todo lo aleatorio usa semilla fija (\codigo{random_state=0} en los modelos, semilla 2026 en el
Monte Carlo), así que correr de nuevo da exactamente lo mismo. La auditoría del 29 de septiembre volvió a correr el Radar
desde cero y las 8 tablas salieron idénticas.

\section{El entorno de cómputo}
\begin{itemize}
\item Python 3.11 en un entorno propio (\codigo{.venv}); todas las versiones fijas en \codigo{requirements.txt}.
\item \textbf{PySpark 3.5.6 con Java 17.} La máquina tenía Java 8, y Spark 3.5 necesita Java 17. Se instaló Java 17 sin
  cambiar el de Windows: \codigo{entorno.py} lo apunta solo dentro del proceso del proyecto.
\item \textbf{winutils y hadoop.dll 3.3.6}: sin ellos, Spark no puede escribir archivos en Windows.
\item \textbf{Ruta corta de Windows}: la carpeta del proyecto tiene espacios (``PP 5to semestre'') y los scripts de Spark
  la cortaban; se usa el nombre corto que Windows asigna sin espacios.
\end{itemize}
Alternativa descartada: Spark dentro de Docker. Funciona, pero es más pesado de arrancar y de explicar en el coloquio.
